\begin{table}[htbp]
\centering
\caption{Diagnóstico tipo Hausman: DR-DiD versus S\&X aproximado}
\label{tab:hausman_composicao}
\small
\begin{tabular}{lccccc}
\toprule
Série-disciplina & DR-DiD & S\&X aprox. & Dif. & \(p\)-valor & MDE 80\% \\
\midrule
9º Ano EF -- LP & 0,510 & 0,502 & -0,008 & 0,9630 & 0,198 \\
9º Ano EF -- Matemática & 0,601 & 0,603 & +0,002 & 0,9976 & 0,219 \\
3ª Série EM -- LP & 0,571 & 0,601 & +0,031 & 0,9000 & 0,248 \\
3ª Série EM -- Matemática & 0,627 & 0,675 & +0,048 & 0,9057 & 0,243 \\
\bottomrule
\end{tabular}
\vspace{0.2cm}
\begin{flushleft}
\footnotesize Nota: os valores de DR-DiD e S\&X correspondem ao efeito médio pós-tratamento, calculado como a média de \(k=0\) a \(k=4\). O teste usa variância conservadora igual à soma das variâncias marginais. MDE 80\% é a diferença mínima detectável, em desvios-padrão, com 80\% de potência e teste bicaudal de 5\%, calculada a partir da raiz da média dos erros-padrão conservadores ao quadrado nas células comparáveis. A coluna S\&X aprox. exclui células sem suporte comum suficiente e não usa fallback binário para DR-DiD.
\end{flushleft}
\end{table}
